%=====================================================================
% FRONT MATTER
%=====================================================================
\pagenumbering{roman}
% Front-matter figures sit outside any chapter, so give them their own prefix
% rather than letting them inherit a misleading chapter number.
\structuralfigures{F.}
\thispagestyle{empty}

\begin{center}
\vspace*{1.0cm}

{\color{secondary}\rule{0.85\textwidth}{2pt}}\\[6pt]
{\large\textsc{\color{secondary}Bachelor of Science in Information Technology}}\\[18pt]

{\Huge\bfseries\color{primary}MACHINE LEARNING}\\[10pt]
{\LARGE\bfseries\color{primary}From Hypotheses to Decisions}\\[16pt]

{\color{secondary}\rule{0.85\textwidth}{2pt}}\\[16pt]

{\Large\textsc{Unified Lecture Handout}}\\[6pt]
{\large\color{gray!70!black}Complete Course in One Volume $\bullet$ Fall 2026}\\[24pt]

% --- The one application domain, announced on the title page ---
\begin{tikzpicture}
  \node[rounded corners=4pt, draw=primary, line width=1.1pt, fill=primary!6,
        text=primary, align=center, font=\small\bfseries,
        minimum width=10.2cm, minimum height=1.5cm]
       {\faIcon{tasks}\\[2pt]Every example drawn from IT project management};
\end{tikzpicture}\\[10pt]
{\small\itshape\color{gray!70!black}Effort and schedules, tasks and sprints, tickets and
releases, risk and technical debt.}

\vspace{1.2cm}

\begin{tcolorbox}[enhanced, width=0.86\textwidth, colback=lightbg,
  colframe=primary, boxrule=0pt, leftrule=4pt, arc=2pt, top=6pt, bottom=6pt]
\small
\textbf{\color{primary}What this volume covers.} The whole of the HEC course
\emph{Machine Learning} --- every topic in the HEC 2023 course outline, and every
learning outcome that HEC 2025 attaches to the course --- in one resequenced
volume, taught from first principles to working code. It is the full-semester
companion to the short Summer 2026 handout, which covered a subset of these
topics for an eight-week session.
\end{tcolorbox}

\vspace{1.0cm}
{\small\color{gray!60!black}Six Parts $\bullet$ Twenty-Two Chapters $\bullet$ Six Appendices}
\vspace{1cm}
\end{center}

\newpage

%=====================================================================
\chapter*{How to Use This Handout}
\addcontentsline{toc}{chapter}{How to Use This Handout}
\markboth{How to Use This Handout}{}

This is not a reference manual to be searched; it is a \term{path} to be walked.
The chapters are ordered so that every idea rests on ideas you have already met.
Reading out of order is possible, but each chapter tells you honestly what it
assumes.

\section*{The shape of the book}

The course moves through six parts, and the movement is deliberate:

\begin{tight}
  \item \textbf{Part I --- Foundations.} What learning from data means, the
        mathematics it runs on, the workflow and toolkit every later chapter
        uses, and the oldest precise account of learning: concept learning as a
        search through a space of hypotheses.
  \item \textbf{Part II --- Linear Models and Generalization.} Linear and
        logistic regression, then the two ideas the rest of the course depends
        on: why a model that fits its training data can still fail, and how to
        measure a model honestly.
  \item \textbf{Part III --- Supervised Learning Algorithms.} Decision trees,
        Bayesian learning, nearest neighbours, support vector machines, neural
        networks, and ensembles that combine them.
  \item \textbf{Part IV --- Learning Without Full Labels.} Clustering,
        dimensionality reduction, self-organizing maps, and semi-supervised
        learning with the EM algorithm.
  \item \textbf{Part V --- Sequences, Decisions and Reinforcement.} Hidden
        Markov models, Markov decision processes, and agents that learn by
        trial and error.
  \item \textbf{Part VI --- Responsible Practice.} The ethics of deploying a
        learned model, and the end-to-end project in which you build one.
\end{tight}

\begin{conceptbox}[Why Evaluation Comes Before the Algorithms]
Part II ends with two chapters about \emph{judging} models rather than building
them: \chref{generalization} and \chref{evaluation}. They come before decision
trees, SVMs and neural networks on purpose. Every algorithm in Part III is
introduced with a claim that it works; without the tools of Part II you have no
way to check that claim, and no way to choose between two algorithms that both
seem to work. Students who learn the algorithms first and evaluation last tend
to report training accuracy for the rest of their careers.
\end{conceptbox}

\section*{The furniture in every chapter}

Each chapter is built from the same parts, always in the same order, so you
always know where you are:

\rowcolors{2}{white}{lightbg}
\begin{longtable}{@{}L{4.6cm}L{10.6cm}@{}}
\toprule
\rowcolor{primary!15}
\textbf{Element} & \textbf{What it is for} \\
\midrule
\endhead
\faIcon{bullseye}~\textbf{Learning Outcomes} & What you will be able to \emph{do} after the chapter. Written as actions, not topics. \\
\faIcon{link}~\textbf{Before You Start} & The specific earlier chapters this one leans on. \\
\faIcon{tags}~\textbf{Key Terms} & Vocabulary introduced here. All terms also appear in the glossary (Appendix~D). \\
\faIcon{book}~\textbf{Definition} & A precise statement of a term. Learn these exactly. \\
\faIcon{lightbulb}~\textbf{Concept} & The intuition behind a definition --- why it is built this way. \\
\faIcon{calculator}~\textbf{Worked Example} & A complete calculation with real numbers. Work through it with pen and paper; examination numericals follow the same pattern. \\
\faIcon{exclamation-circle}~\textbf{Common Pitfalls} & Mistakes that are made every semester. Read these twice. \\
\faIcon{flask}~\textbf{Try It Yourself} & A short Python lab you can run immediately. Every lab is also in the \texttt{code/} folder as a runnable file. \\
\faIcon{question-circle}~\textbf{Checkpoint} & Three quick self-tests. Answers in Appendix~C. \\
\faIcon{clipboard-check}~\textbf{Chapter Summary} & The chapter compressed to what must be retained. \\
\faIcon{pen-fancy}~\textbf{Review Questions} & Multiple-choice, short-answer and applied questions for assessment practice. \\
\bottomrule
\end{longtable}

\section*{How to study with it}

\begin{tightnum}
  \item \textbf{Read the outcomes first.} They tell you what to look for.
  \item \textbf{Do not skip the worked examples.} Reading a calculation is not
        the same as being able to perform one. Cover the answer and try it.
  \item \textbf{Answer the checkpoint before moving on.} If you cannot, the
        next chapter will be harder than it needs to be.
  \item \textbf{Run the labs.} Fifteen minutes of running code teaches more
        than an hour of reading about it. Then change one number and predict
        what will happen before you run it again.
  \item \textbf{Connect each algorithm to a project decision.} For every chapter,
        say which project-management question it could answer --- how long, how
        risky, which queue, when to refactor. If you can, you understand it.
\end{tightnum}

\begin{alertbox}[A Note on Notation and Rigour]
This handout uses mathematics where mathematics clarifies --- the gradient
descent update, Bayes' theorem, the margin of an SVM, the Bellman equation ---
but it proves nothing that is not needed to use the result correctly. Every
formula is accompanied by an explanation in words and, wherever possible, a
numeric example you can reproduce by hand. If a formula looks intimidating, read
the words around it first and return to the symbols afterwards.
\end{alertbox}

\newpage

%=====================================================================
\chapter*{Course Specification}
\addcontentsline{toc}{chapter}{Course Specification}
\markboth{Course Specification}{}

This page is the course file. It states what HEC requires of a course called
\emph{Machine Learning}, and shows where in this handout each requirement is
met.

\section*{Course particulars}

\rowcolors{2}{white}{lightbg}
\begin{longtable}{@{}L{4.2cm}L{11.0cm}@{}}
\toprule
\rowcolor{primary!15}
\textbf{Item} & \textbf{Value} \\
\midrule
\endhead
Course title & Machine Learning \\
Programme & BS Information Technology (domain elective); the same course is a
BSAI domain core and a BSDS domain elective under HEC 2023, and an elective in
the Data Science, AI and Software Engineering specializations of the HEC 2025
BS Computer Science \\
Credit hours & 3 (2--3): two credit hours of lectures and one credit hour of
laboratory, which is three contact hours of lab per week \\
Prerequisite & Artificial Intelligence. Probability and statistics, and
programming in Python, are assumed (\chref{math} and \chref{workflow} revise what
is needed) \\
Duration & Sixteen teaching weeks, with a midterm after Week~8 and a final
examination after Week~16 (Appendix~F) \\
Course introduction & The HEC aims, verbatim: ``a) Present the basic machine
learning concepts; b)~Present a range of machine learning algorithms along with
their strengths and weaknesses; c)~Apply machine learning algorithms to solve
problems of moderate complexity.'' \\
\bottomrule
\end{longtable}

\section*{Course learning outcomes}

HEC 2023 gives the course five learning outcomes. HEC 2025 publishes a different
list for each specialization that carries the course; together they add
evaluation, bias--variance reasoning, modern toolkits and ethics. The seven
outcomes below are the union, and each is traced to the chapters that deliver it.

\rowcolors{2}{white}{lightbg}
\begin{longtable}{@{}L{1.3cm}L{7.6cm}L{2.1cm}L{3.8cm}@{}}
\toprule
\rowcolor{primary!15}
\textbf{CLO} & \textbf{By the end of the course the student will be able to} &
\textbf{Bloom} & \textbf{Chapters} \\
\midrule
\endhead
CLO-1 & Explain core machine learning concepts, paradigms and theories, and
select an appropriate model for a task & C2 Understand & 1--4, 7 \\
CLO-2 & Apply supervised learning techniques to classification and regression
problems of moderate complexity & C3 Apply & 5, 6, 9--14 \\
CLO-3 & Apply unsupervised and semi-supervised techniques --- clustering,
dimensionality reduction, self-organizing maps, EM --- to real data & C3 Apply &
15--17 \\
CLO-4 & Apply reinforcement learning algorithms to environments with sequential,
uncertain dynamics & C3 Apply & 18--20 \\
CLO-5 & Analyze the performance of machine learning models using evaluation
metrics, validation and bias--variance trade-offs & C4 Analyze & 7, 8, and every
lab \\
CLO-6 & Develop and test an end-to-end machine learning solution of reasonable
size using modern toolkits & C6 Create & 3, 22, and every lab \\
CLO-7 & Discuss the ethical considerations of applying machine learning to
societal problems & A3 Valuing & 21, and the fairness check of 22 \\
\bottomrule
\end{longtable}

\section*{Coverage of the HEC 2023 course outline}

Every topic in the published outline, in the order HEC lists it, with the
chapter that teaches it. Nothing is omitted.

\rowcolors{2}{white}{lightbg}
\begin{longtable}{@{}L{6.4cm}L{2.0cm}L{6.6cm}@{}}
\toprule
\rowcolor{primary!15}
\textbf{HEC 2023 outline topic} & \textbf{Chapter} & \textbf{Where exactly} \\
\midrule
\endhead
Introduction to machine learning & 1, 2, 3 & Definition, paradigms, the design
of a learning system; the mathematics and the workflow it needs \\
Concept learning: general-to-specific ordering of hypotheses & 4 & The
more-general-than relation; \textsc{Find-S} \\
Version spaces algorithm & 4 & Version spaces; \textsc{List-Then-Eliminate} \\
Candidate elimination algorithm & 4 & Full trace on the release-decision data;
inductive bias \\
Supervised learning: decision trees & 9 & ID3, entropy and information gain,
gain ratio, Gini \\
Naive Bayes & 10 & Bayes' theorem, MAP hypotheses, naive Bayes with smoothing \\
Artificial neural networks & 13 & Perceptron, delta rule, multilayer networks,
backpropagation \\
Support vector machines & 12 & Maximum margin, soft margin, kernels \\
Overfitting, noisy data, and pruning & 7, 9 & Bias--variance and
regularisation (7); pruning trees against noise (9) \\
Measuring classifier accuracy & 8 & Confusion matrix, precision, recall, ROC,
confidence intervals for error, comparing two classifiers \\
Linear and logistic regression & 5, 6 & Least squares and gradient descent;
the sigmoid and cross-entropy \\
Unsupervised learning: hierarchical agglomerative clustering & 15 & Linkage
criteria, dendrograms \\
k-means partitional clustering & 15 & Lloyd's algorithm, choosing $k$ \\
Self-organizing maps (SOM) & 16 & Competitive learning, neighbourhoods, the
U-matrix \\
k-nearest-neighbour algorithm & 11 & Distance, choice of $k$, weighting,
the curse of dimensionality \\
Semi-supervised learning with EM using labeled and unlabeled data & 17 &
Gaussian mixtures, the EM algorithm, EM over labelled plus unlabelled data \\
Reinforcement learning: hidden Markov models & 18 & Forward algorithm,
Viterbi, Baum--Welch as EM \\
Monte Carlo inference & 20 & Estimation by sampling; Monte Carlo value
estimation and control \\
Exploration vs.\ exploitation trade-off & 20 & Multi-armed bandits,
$\varepsilon$-greedy, UCB \\
Markov decision processes & 19 & Bellman equations, value and policy iteration \\
Ensemble learning: committees of multiple hypotheses; bagging, boosting & 14 &
Voting, bagging and random forests, AdaBoost and gradient boosting \\
\bottomrule
\end{longtable}

\begin{conceptbox}[Three Places Where This Handout Departs from the Outline's Grouping]
The HEC outline is a list of topics, not a teaching order, and three of its
groupings would mislead a student if followed literally.
\begin{tight}
  \item \textbf{$k$-nearest-neighbour} is listed under \emph{unsupervised}
        learning. It is a supervised method --- it needs labels to vote with ---
        so it is taught in Part III with the other classifiers (\chref{knn}).
  \item \textbf{Hidden Markov models} are listed under \emph{reinforcement}
        learning. An HMM makes no decisions and receives no reward; it is a
        model of sequences. It is taught at the start of Part V because its
        Markov machinery is exactly what Markov decision processes build on
        (\chref{hmm}).
  \item \textbf{Overfitting and measuring accuracy} are listed after the
        supervised algorithms. They are taught before them (Part II), for the
        reason given on the previous page.
\end{tight}
Every topic is still taught; only the order changes.
\end{conceptbox}

\section*{Coverage of the HEC 2025 learning outcomes}

HEC 2025 publishes learning outcomes but no outline. The outcomes it attaches to
\emph{Machine Learning} in each specialization are all met:

\rowcolors{2}{white}{lightbg}
\begin{longtable}{@{}L{3.1cm}L{8.2cm}L{3.7cm}@{}}
\toprule
\rowcolor{primary!15}
\textbf{Specialization} & \textbf{Distinctive outcomes it adds} & \textbf{Met in} \\
\midrule
\endhead
Software Engineering & Analyze performance using evaluation metrics (C4);
implement with modern toolkits (C3); discuss ethical considerations (A3) &
Ch 8; Ch 3 and labs; Ch 21 \\
Data Science & Apply ML to classification problems and to real-world problems
(C3) & Chs 6, 9--14; project (Ch 22) \\
Artificial Intelligence & Differentiate ML paradigms and select appropriate
models; implement, train and evaluate SVMs, decision trees and $k$-means;
analyse performance with metrics & Chs 1, 7; Chs 9, 12, 15; Ch 8 \\
Robotics & Supervised, unsupervised and reinforcement learning; regression, SVM,
PCA and clustering; bias--variance trade-offs & Chs 1, 5, 12, 15, 16, 19--20;
Ch 7 \\
Network Infrastructure and Cloud Computing & Design and develop end-to-end ML
solutions for real-world datasets & Ch 22 and the labs \\
\bottomrule
\end{longtable}

\begin{alertbox}[One HEC 2025 List Belongs to a Different Course]
The Internet of Things specialization lists, under the name \emph{Machine
Learning}, outcomes about CNNs, RNNs and transformers and about deploying models
on devices. Those are the outcomes of a \emph{Deep Learning} course. This handout
introduces neural networks thoroughly (\chref{ann}) and signposts the deep
architectures, but does not teach them; a department offering this course in the
IoT specialization should pair it with Deep Learning rather than stretch it.
\end{alertbox}

\section*{What this course deliberately leaves to other courses}

Machine Learning overlaps heavily with Data Mining and with Introduction to Data
Science. Following the split recommended in the department's HEC curriculum map,
this course owns \emph{prediction} --- generalization, model selection,
evaluation, ensembles and regularisation --- and leaves \emph{discovery} to Data
Mining: association rules and frequent patterns, large-scale preprocessing,
outlier mining and graph mining are not taught here. Deep architectures (CNNs,
RNNs, transformers) belong to Deep Learning, and deployment and monitoring to
MLOps.

\section*{Assessment}

\rowcolors{2}{white}{lightbg}
\begin{longtable}{@{}L{4.4cm}L{1.6cm}L{9.0cm}@{}}
\toprule
\rowcolor{primary!15}
\textbf{Instrument} & \textbf{Weight} & \textbf{What it assesses} \\
\midrule
\endhead
Quizzes (three) & 10\% & CLO-1, CLO-2: definitions and one short numerical each \\
Assignments (two) & 10\% & CLO-2, CLO-5: a by-hand numerical and a coded
experiment each \\
Project & 15\% & CLO-5, CLO-6, CLO-7: the end-to-end project of \chref{project} \\
Midterm examination & 25\% & Chapters 1--8 \\
Final examination & 40\% & The whole course, weighted towards Chapters 9--21 \\
\bottomrule
\end{longtable}

{\small\itshape These weights are a suggestion. Replace them with your
department's scheme; the mapping from instruments to outcomes is what a course
file actually needs to show.}

\Needspace{10\baselineskip}
\section*{Recommended texts}

\textbf{Set by HEC.}
\begin{tightnum}
  \item Tom M.\ Mitchell, \emph{Machine Learning}, McGraw-Hill, 1997. The source
        of the concept-learning, decision-tree, Bayesian-learning and
        instance-based chapters; its worked examples are reproduced here.
  \item Kevin P.\ Murphy, \emph{Machine Learning: A Probabilistic Perspective},
        MIT Press, 2012. The reference for EM, HMMs and the probabilistic view.
\end{tightnum}

\textbf{Recommended alongside.}
\begin{tightnum}
  \setcounter{enumi}{2}
  \item Aur\'elien G\'eron, \emph{Hands-On Machine Learning with Scikit-Learn,
        Keras and TensorFlow}, 3rd ed., O'Reilly, 2022. For the labs.
  \item Gareth James, Daniela Witten, Trevor Hastie, Robert Tibshirani and
        Jonathan Taylor, \emph{An Introduction to Statistical Learning with
        Applications in Python}, Springer, 2023. Free online; the clearest
        treatment of the bias--variance trade-off.
  \item Richard S.\ Sutton and Andrew G.\ Barto, \emph{Reinforcement Learning:
        An Introduction}, 2nd ed., MIT Press, 2018. Free online; the source for
        Chapters 19 and 20.
\end{tightnum}

\newpage

%=====================================================================
\chapter*{The Course at a Glance}
\addcontentsline{toc}{chapter}{The Course at a Glance}
\markboth{The Course at a Glance}{}

\dsfigh{%
\resizebox{\textwidth}{!}{%
\begin{tikzpicture}[
  part/.style={rounded corners=3pt, fill=#1, text=white, font=\bfseries\footnotesize,
               minimum width=2.7cm, minimum height=0.9cm, align=center},
  ch/.style={rounded corners=2pt, draw=#1, line width=0.7pt, fill=#1!7, text=#1!60!black,
             font=\scriptsize, minimum width=2.7cm, minimum height=0.52cm, align=center},
  arr/.style={-{Stealth[length=2mm]}, primary!45, line width=0.8pt}
]
\node[part=primary]   (p1) at (0,0)    {Part I\\Foundations};
\node[part=secondary] (p2) at (3.1,0)  {Part II\\Linear Models};
\node[part=greenhl]   (p3) at (6.2,0)  {Part III\\Supervised};
\node[part=purplehl]  (p4) at (9.3,0)  {Part IV\\Without Labels};
\node[part=tealhl]    (p5) at (12.4,0) {Part V\\Decisions \& RL};
\node[part=goldhl]    (p6) at (15.5,0) {Part VI\\Practice};

\foreach \a/\b in {p1/p2,p2/p3,p3/p4,p4/p5,p5/p6}{\draw[arr] (\a) -- (\b);}

\foreach \i/\t in {1/1. Introduction, 2/2. Mathematics, 3/3. Workflow,
                   4/4. Concept Learning}{
  \pgfmathsetmacro{\yy}{-0.85-(\i-1)*0.62}
  \node[ch=primary] at (0,\yy) {\t};}

\foreach \i/\t in {1/5. Linear Regression, 2/6. Logistic Regression,
                   3/7. Generalization, 4/8. Evaluation}{
  \pgfmathsetmacro{\yy}{-0.85-(\i-1)*0.62}
  \node[ch=secondary] at (3.1,\yy) {\t};}

\foreach \i/\t in {1/9. Decision Trees, 2/10. Bayesian Learning,
                   3/11. Nearest Neighbours, 4/12. SVMs,
                   5/13. Neural Networks, 6/14. Ensembles}{
  \pgfmathsetmacro{\yy}{-0.85-(\i-1)*0.62}
  \node[ch=greenhl] at (6.2,\yy) {\t};}

\foreach \i/\t in {1/15. Clustering, 2/16. PCA and SOM, 3/17. Semi-supervised}{
  \pgfmathsetmacro{\yy}{-0.85-(\i-1)*0.62}
  \node[ch=purplehl] at (9.3,\yy) {\t};}

\foreach \i/\t in {1/18. Hidden Markov, 2/19. MDPs, 3/20. Reinforcement}{
  \pgfmathsetmacro{\yy}{-0.85-(\i-1)*0.62}
  \node[ch=tealhl] at (12.4,\yy) {\t};}

\foreach \i/\t in {1/21. Ethics, 2/22. Project}{
  \pgfmathsetmacro{\yy}{-0.85-(\i-1)*0.62}
  \node[ch=goldhl] at (15.5,\yy) {\t};}

% the two hinge chapters of Part II
\foreach \yy/\t in {-2.09/7. Generalization, -2.71/8. Evaluation}{
  \node[rounded corners=2pt, draw=secondary, line width=1.6pt, fill=secondary!18,
        text=secondary!50!black, font=\scriptsize\bfseries, minimum width=2.7cm,
        minimum height=0.52cm] at (3.1,\yy) {\t};}

% ---- midterm marker
\draw[accent, dashed, line width=0.9pt] (4.65,0.7) -- (4.65,-4.35);
\node[font=\tiny\bfseries, text=accent, fill=white, inner sep=1.5pt] at (4.65,-4.35)
     {midterm};

% ---- the one running example underneath everything
\node[rounded corners=3pt, fill=primary!8, draw=primary!30,
      minimum width=18.3cm, minimum height=0.62cm,
      font=\scriptsize\bfseries, text=primary] at (7.75,-5.0)
      {\faIcon{tasks}~~Running example in every chapter: IT project management
       --- effort, schedules, tickets, releases, risk and technical debt};
\end{tikzpicture}}}%
{Course map. The six parts run left to right in the order they are taught; the chapters
in each column are the order within that part. The two outlined chapters of Part II are
the hinge of the course, and the dashed line marks the midterm.}{coursemap}

\section*{The Running Example: IT Project Management}

Every example in this handout comes from one setting: a software organisation
delivering IT projects. Keeping to one setting keeps the course simple --- the
vocabulary of tasks, sprints, tickets and releases is learned once, and each new
algorithm is met as a new answer to a familiar kind of question.

\rowcolors{2}{white}{lightbg}
\begin{longtable}{@{}L{8.4cm}L{6.6cm}@{}}
\toprule
\rowcolor{primary!15}
\textbf{Recurring example} & \textbf{Chapters} \\
\midrule
\endhead
Story points and the hours stories really take & 1, 5 \\
Tasks and projects that finish late & 1, 3, 6, 11--14, 16, 22 \\
Release approvals and sprint goals & 4, 9, 10 \\
Code changes, defects and releases & 2, 8, 10 \\
Tickets: their text, type, queue and resolution time & 6, 8, 10, 15, 17 \\
A project's hidden weekly health & 18, 20 \\
Technical debt: ship features or refactor? & 19, 20 \\
Screening candidates for developer roles & 21 \\
\bottomrule
\end{longtable}

\begin{conceptbox}[What Makes Project Data Hard]
Project data is \emph{small} --- a team completes a few hundred tasks a year --- so
overfitting (\chref{generalization}) is a constant danger. It is \emph{noisy and
human-reported}: estimates, statuses and root causes are typed in by busy people.
Its interesting events are \emph{rare}: most tasks finish, most changes are clean,
so accuracy misleads (\chref{evaluation}). And a prediction \emph{changes the
behaviour} of the team that hears it, so a model of last year's projects describes
a world the model itself will alter. Every chapter meets at least one of these.
\end{conceptbox}

\noindent Two of the set text's classic tables --- Mitchell's \emph{EnjoySport}
(\chref{concept}) and \emph{PlayTennis} (\chref{trees}, \chref{bayes}) --- appear here
with their attributes renamed as release and sprint decisions, so every number can
still be checked against the book.

\section*{Notation Used in This Handout}

\rowcolors{2}{white}{defbg}
\begin{longtable}{@{}L{2.8cm}L{12.4cm}@{}}
\toprule
\rowcolor{primary!15}
\textbf{Symbol} & \textbf{Meaning} \\
\midrule
\endhead
$n$ & Number of training examples (rows) \\
$d$ & Number of features (columns, attributes) \\
$\mathbf{x}_i$ & The $i$-th example, a vector of $d$ feature values; bold lower case is always a vector \\
$x_{ij}$ & The value of feature $j$ in example $i$ \\
$y_i$ & The true target (label) of example $i$ \\
$\hat{y}_i$ & The model's \emph{prediction} of $y_i$; the hat always means ``estimated'' \\
$X$ & The feature matrix, $n$ rows by $d$ columns; capitals are matrices \\
$\mathbf{w},\ b$ & Weights and bias (intercept) of a linear model \\
$\theta$ & All parameters of a model, collected together \\
$h,\ H$ & A hypothesis, and the hypothesis space it is drawn from (\chref{concept}) \\
$L(\theta)$ & Loss --- how wrong the model with parameters $\theta$ is on the data \\
$\nabla L$ & Gradient: the vector of partial derivatives of $L$ \\
$\alpha$ or $\eta$ & Learning rate (step size) \\
$\lambda$ & Regularisation strength \\
$P(A \mid B)$ & Probability of $A$ given $B$ \\
$\mathbb{E}[X]$ & Expected (average) value of $X$ \\
$s,\ a,\ r$ & State, action and reward, in Part V \\
$\gamma$ & Discount factor for future rewards, in Part V \\
$\pi$ & A policy (Part V); also the initial-state distribution of an HMM (\chref{hmm}) \\
\bottomrule
\end{longtable}

\begin{alertbox}[Two Symbols Students Confuse Every Year]
$y$ versus $\hat{y}$: $y$ is the \emph{truth}, $\hat{y}$ is the \emph{guess}.
Every loss function in this book is some way of summarising the gaps between
them.\\[3pt]
\emph{Training} error versus \emph{true} error: the first is measured on the data
the model learned from, the second is what happens on data it has never seen.
The whole of Part II is about the gap between them.
\end{alertbox}

\section*{Prerequisite Self-Check}

You are ready for Part I if you can answer these. If several give you trouble,
work through \chref{math} slowly and keep Appendix~B beside you.

\begin{tightnum}
  \item Compute the dot product of $(1, 2, 3)$ and $(4, 0, -1)$.
  \item What is the derivative of $f(w) = (w - 3)^2$, and where is $f$ smallest?
  \item A test is positive for 90\% of sick people and 10\% of healthy ones; 1\%
        of people are sick. Roughly what fraction of positive tests are sick
        people?
  \item What is $\log_2 8$? What is $\log_2 1$?
  \item In a table of network connections, what does one \emph{row} usually
        represent? What does one \emph{column} represent?
  \item Write a Python loop that prints the squares of 1 to 5.
\end{tightnum}

\smallskip
\textit{Answers: (1)~$4+0-3=1$\quad (2)~$2(w-3)$, smallest at $w=3$\quad
(3)~about 8\% --- $0.009/(0.009+0.099)$\quad (4)~3; 0\quad
(5)~one row = one connection (an example); one column = one attribute (a
feature)\quad (6)~e.g.\ \texttt{for i in range(1,6): print(i**2)}}

\newpage
\tableofcontents
\newpage
\listoffigures
\addcontentsline{toc}{chapter}{List of Figures}
% No List of Tables: tables in this handout are inline reference material rather
% than numbered exhibits, so they are not captioned and a LoT would be empty.
\newpage

\pagenumbering{arabic}
\setcounter{page}{1}
